% ==============================================================================
% PREAMBLE - ISTITUZIONI DI ALGEBRA (PROF. GIOVANNI CUTOLO)
% Corso di Laurea in Informatica - Università degli Studi di Napoli Federico II
% Curatore: Fagghino
% ==============================================================================

\usepackage[utf8]{inputenc}
\usepackage[T1]{fontenc}
\usepackage[italian]{babel}

% Geometria pagina e margini
\usepackage[a4paper,top=2.5cm,bottom=2.5cm,left=2.5cm,right=2.5cm,headheight=14.5pt]{geometry}

% Matematica e simboli
\usepackage{amsmath,amssymb,amsthm,mathtools,amsfonts,mathrsfs}
\usepackage{bm}

% Grafica e diagrammi
\usepackage{tikz}
\usepackage{tikz-cd}
\usetikzlibrary{shapes,arrows,positioning,calc,matrix,fit,backgrounds}

% Box decorati per definizioni, teoremi ed osservazioni
\usepackage[most]{tcolorbox}

% Gestione elenchi
\usepackage{enumitem}
\setlist{nosep}

% Intestazioni e piè di pagina
\usepackage{fancyhdr}
\pagestyle{fancy}
\fancyhf{}
\fancyhead[LE,RO]{\thepage}
\fancyhead[RE]{\nouppercase{\leftmark}}
\fancyhead[LO]{\nouppercase{\rightmark}}
\renewcommand{\headrulewidth}{0.4pt}

% Collegamenti ipertestuali e riferimenti
\usepackage{xcolor}
\definecolor{UniNaBlue}{RGB}{0, 48, 135}
\definecolor{CutoloRed}{RGB}{160, 20, 20}
\definecolor{DefGreen}{RGB}{20, 110, 60}
\definecolor{ExOrange}{RGB}{180, 80, 0}
\definecolor{NoteGray}{RGB}{80, 80, 90}

\usepackage{hyperref}
\hypersetup{
    colorlinks=true,
    linkcolor=UniNaBlue,
    citecolor=CutoloRed,
    urlcolor=UniNaBlue,
    pdftitle={Istituzioni di Algebra - Appunti Ufficiali Prof. Cutolo},
    pdfauthor={Fagghino},
    pdfsubject={Corso di Algebra - Università degli Studi di Napoli Federico II}
}

\usepackage{cleveref}

% ==============================================================================
% MACRO MATEMATICHE E NOTAZIONI UFFICIALI CUTOLO
% ==============================================================================

% Insiemi numerici standard
\newcommand{\N}{\mathbb{N}}
\newcommand{\Z}{\mathbb{Z}}
\newcommand{\Q}{\mathbb{Q}}
\newcommand{\R}{\mathbb{R}}
\newcommand{\C}{\mathbb{C}}
\newcommand{\K}{\mathbb{K}}
\newcommand{\F}{\mathbb{F}}

% Insieme delle parti
\renewcommand{\P}{\mathscr{P}}
\newcommand{\Pfin}{\mathscr{P}_{\text{fin}}}

% Relazioni e connettivi speciali
\newcommand{\vuoto}{\varnothing}
\newcommand{\smin}{\setminus}
\newcommand{\implica}{\Rightarrow}
\newcommand{\Implica}{\Longrightarrow}
\newcommand{\shiff}{\Leftrightarrow}
\newcommand{\Shiff}{\Longleftrightarrow}
\newcommand{\sseq}{\subseteq}
\newcommand{\ds}{\mathbin{\scriptstyle\triangle}} % Differenza simmetrica
\newcommand{\xor}{\mathbin{\mathsf{XOR}}}
\newcommand{\nor}{\mathbin{\mathsf{NOR}}}
\newcommand{\nand}{\mathbin{\mathsf{NAND}}}

% Notazione specifica del Prof. Cutolo per immagini e controimmagini
% Freccia verso destra per l'immagine f(X)
\newcommand{\vecf}{\vec{f}}
% Freccia verso sinistra per la controimmagine / anti-immagine (NON f^-1!)
\newcommand{\antivec}[2]{#1^{\raise #2pt\hbox{$\!\!\scriptstyle\leftarrow\!\!$}}}
\newcommand{\antivecf}{\antivec{f}{3}}
\newcommand{\antivecg}{\antivec{g}{2}}

% Operatori e insiemi di strutture
\DeclareMathOperator{\id}{id}
\DeclareMathOperator{\im}{im}
\DeclareMathOperator{\Ker}{Ker}
\DeclareMathOperator{\Sym}{Sym}
\DeclareMathOperator{\Aut}{Aut}
\DeclareMathOperator{\Hom}{Hom}
\DeclareMathOperator{\End}{End}
\DeclareMathOperator{\Eq}{Eq}
\DeclareMathOperator{\Part}{Part}
\DeclareMathOperator{\Rel}{Rel}
\DeclareMathOperator{\Corr}{Corr}
\DeclareMathOperator{\Map}{Map}
\DeclareMathOperator{\car}{char}
\DeclareMathOperator{\mcd}{MCD}
\DeclareMathOperator{\mcm}{mcm}
\DeclareMathOperator{\minor}{Minor}
\DeclareMathOperator{\maggior}{Maggior}
\DeclareMathOperator{\degpoly}{deg}

% Generatori e classi resto
\newcommand{\gen}[1]{\langle #1 \rangle}
\newcommand{\classe}[1]{[#1]}
\newcommand{\modulo}[1]{\ (\mathrm{mod}\ #1)}

% ==============================================================================
% AMBIENTI TCOLORBOX (Definizioni, Teoremi, Proposizioni, Esercizi)
% ==============================================================================

% Stile per Definizioni
\newtcbtheorem[number within=chapter]{definizione}{Definizione}{
    enhanced,
    colback=DefGreen!5!white,
    colframe=DefGreen!80!black,
    fonttitle=\bfseries,
    separator sign dash,
    boxrule=0.8pt,
    arc=3pt,
    top=6pt,bottom=6pt,left=8pt,right=8pt
}{def}

% Stile per Teoremi
\newtcbtheorem[use counter from=definizione]{teorema}{Teorema}{
    enhanced,
    colback=UniNaBlue!5!white,
    colframe=UniNaBlue!90!black,
    fonttitle=\bfseries,
    separator sign dash,
    boxrule=0.8pt,
    arc=3pt,
    top=6pt,bottom=6pt,left=8pt,right=8pt
}{teo}

% Stile per Proposizioni / Lemmi
\newtcbtheorem[use counter from=definizione]{proposizione}{Proposizione}{
    enhanced,
    colback=UniNaBlue!3!white,
    colframe=UniNaBlue!60!black,
    fonttitle=\bfseries,
    separator sign dash,
    boxrule=0.6pt,
    arc=2pt,
    top=5pt,bottom=5pt,left=7pt,right=7pt
}{prop}

% Stile per Esercizi e Problemi Svolti
\newtcbtheorem[number within=chapter]{esercizio}{Esercizio d'Esame}{
    enhanced,
    colback=ExOrange!4!white,
    colframe=ExOrange!85!black,
    fonttitle=\bfseries,
    separator sign dash,
    boxrule=0.8pt,
    arc=3pt,
    top=6pt,bottom=6pt,left=8pt,right=8pt
}{es}

% Box per i Tranelli dell'Orale e le Osservazioni di Cutolo
\newtcolorbox{avvisoCutolo}[1][]{
    enhanced,
    colback=CutoloRed!5!white,
    colframe=CutoloRed!85!black,
    fonttitle=\bfseries,
    title={\textbf{All'Orale con Cutolo: Attenzione al Rigore!}},
    boxrule=0.8pt,
    arc=3pt,
    top=6pt,bottom=6pt,left=8pt,right=8pt,
    #1
}

% Box per l'intuizione o la panoramica
\newtcolorbox{intuizione}[1][]{
    enhanced,
    colback=gray!5!white,
    colframe=gray!60!black,
    fonttitle=\bfseries,
    title={Intuizione e Quadro Generale},
    boxrule=0.6pt,
    arc=2pt,
    top=5pt,bottom=5pt,left=7pt,right=7pt,
    #1
}
